\section{Fundamentação técnica}

O fluxo do trabalho é composto por ingestão RTMP, transcodificação em múltiplas variantes e disponibilização via HLS. O padrão HLS é descrito no RFC 8216, que orienta a segmentação do vídeo e a composição de playlists adaptativas \cite{rfc8216}. Na prática, a solução utiliza servidores e transcodificadores distintos para isolar o impacto de cada componente no desempenho final.

O FFmpeg é utilizado como transcodificador por sua ampla adoção e flexibilidade em pipelines de áudio e vídeo \cite{ffmpeg-docs}. Já o GStreamer oferece uma abordagem de construção por elementos, com boa capacidade de composição de fluxos multimídia \cite{gstreamer-docs}. No lado da ingestão, Nginx-RTMP fornece uma base enxuta e tradicional para RTMP, enquanto o SRS oferece um servidor especializado em streaming em tempo real \cite{nginx-rtmp-module,srs-docs}.

O projeto está empacotado em containers Docker, o que facilita o isolamento do ambiente e a repetição dos testes \cite{docker-docs}. O backend em Spring Boot coordena parte do ciclo de vida da aplicação e expõe a camada de controle usada pelos scripts de benchmark \cite{spring-boot-docs}. O frontend, por sua vez, atua como superfície de reprodução e observação do comportamento do HLS durante os testes.

\subsection{Leitura correta das métricas}

As métricas de CPU e memória precisam ser lidas sempre em relação ao host documentado, porque o percentual de CPU do container pode ultrapassar 100\% quando há execução concorrente em múltiplas threads. Isso não indica erro por si só; indica a intensidade de uso agregada pelo container. A mesma lógica vale para memória: o consumo observado deve ser comparado com a capacidade total do equipamento e com a carga dos demais serviços em execução.

No conjunto atual, não há uso de GPU no pipeline. Portanto, qualquer ganho ou perda de desempenho observada nos resultados deve ser atribuída ao comportamento do software, ao ambiente Docker e ao hardware do host, e não a aceleração por dispositivos gráficos.